\section{Modelo Conceptual del Portal}

El modelo conceptual del Portal Agro Comercial del Huila se basa en una arquitectura modular e integrada, inspirada en las plataformas tecnológicas descritas por Ojeda-Beltrán (2022) y Wolfert et al. (2017). Este modelo conceptualiza el portal como un ecosistema digital que conecta actores de la cadena de valor agrícola, facilitando transacciones transparentes y reduciendo intermediación.

\subsection{Componentes Principales}

El modelo se estructura en cuatro capas interconectadas:

1. **Capa de Datos**: Integra IoT, big data y geoespaciales para recopilar y analizar información de fincas, clima y mercados.

2. **Capa de Servicios**: Incluye marketplace e-commerce, trazabilidad blockchain y herramientas de agricultura de precisión.

3. **Capa de Usuarios**: Interfaces adaptadas para productores, compradores e intermediarios, con módulos de capacitación.

4. **Capa de Gobernanza**: Políticas de inclusión digital, alianzas y monitoreo de impacto.

\subsection{Flujo de Interacción}

Los productores registran fincas con datos IoT, publican ofertas en el marketplace. Compradores acceden a información trazable, realizando transacciones seguras. El sistema optimiza cadenas de suministro con IA.

\subsection{Integración Tecnológica}

El modelo integra TIC para empoderar agricultores, como apps móviles y RRSS, cerrando brechas digitales mediante capacitación.

Gráfica sugerida: Diagrama conceptual del portal mostrando capas y flujos de interacción.

Este modelo reduce intermediación al democratizar acceso a mercados, promoviendo sostenibilidad en el Huila.